\section{Protocol Hecke-like operators for refinement under uplift (notes)}
\label{app:protocol_hecke_operators}

\AuditTag This appendix records a minimal, fully computable ``operator algebra'' viewpoint on window refinement.
It turns the extension/refinement multiplicities under uplift into a semigroup of explicit linear operators, and highlights a rigid second-order recurrence (a prime-power skeleton) shared with classical Hecke templates.
The results are mathematical bookkeeping statements; they introduce no new axiom and are not used as premises in theorem-level folding proofs.

\subsection{Two-state refinement operators from the golden-mean transition matrix}
\label{subsec:protocol_hecke_two_state}

Let $A=\bigl(\begin{smallmatrix}1&1\\1&0\end{smallmatrix}\bigr)$ be the golden-mean transition matrix (Lemma~\ref{lem:golden_zeta_matrix}).
We regard vectors in $\CC^2$ as functions on the two ``boundary states'' $\{0,1\}$ (the last bit of a prefix).

\begin{definition}[Refinement operators $T_L$]\label{def:refinement_operators_TL}
\MathTag For each $L\in\ZZ_{\ge 0}$ define the linear operator
\[
T_L:\CC^2\to\CC^2,\qquad T_L v:=A^L v.
\]
Equivalently, $(T_L)_{ij}$ counts admissible length-$L$ extensions from state $i$ to state $j$ in the golden-mean transition graph.
\end{definition}

\begin{proposition}[Semigroup and commutativity]\label{prop:TL_semigroup_commute}
\MathTag For all $L,M\ge 0$ one has
\[
T_L\circ T_M=T_{L+M},\qquad T_L T_M=T_M T_L.
\]
\end{proposition}
\begin{proof}
This is the law of matrix powers: $A^{L}A^{M}=A^{L+M}$.
Commutativity follows because all $T_L$ are polynomials in the same matrix $A$.
\end{proof}

\begin{proposition}[Prime-power recurrence (Cayley--Hamilton)]\label{prop:TL_prime_power_recurrence}
\MathTag Let $t:=\mathrm{tr}(A)=1$ and $\delta:=\det(A)=-1$.
Then the family $(T_L)$ satisfies the second-order recurrence
\[
T_{L+1}=t\,T_L-\delta\,T_{L-1}=T_L+T_{L-1}\qquad(L\ge 1),
\]
with initial conditions $T_0=I$ and $T_1=A$.
\end{proposition}
\begin{proof}
By Cayley--Hamilton, $A^2-(\mathrm{tr}A)A+(\det A)I=0$.
Multiplying by $A^{L-1}$ gives $A^{L+1}=(\mathrm{tr}A)A^L-(\det A)A^{L-1}$ for $L\ge 1$, i.e.\ the recurrence for $T_L=A^L$.
\end{proof}

\subsection{Extension multiplicities as operator evaluations}
\label{subsec:ext_counts_operator_eval}

Fix the base window $6$ and let $u=u_1\cdots u_6\in X_6$.
For $m\ge 6$ write $L=m-6$.
Recall
\[
\mathrm{Ext}_m(u):=\{w\in X_m:\ \pi_{m\to 6}(w)=u\}
\]
from Proposition~\ref{prop:prefix_extension_count}.

\begin{proposition}[Operator expression for $|\mathrm{Ext}_m(u)|$]\label{prop:ext_count_operator_formula}
\MathTag Let $\mathbf{1}=(1,1)^{\mathsf{T}}$ and $e_0=(1,0)^{\mathsf{T}}$, $e_1=(0,1)^{\mathsf{T}}$.
Then for $L=m-6$ one has
\[
|\mathrm{Ext}_m(u)|=e_{u_6}^{\mathsf{T}}\,T_L\,\mathbf{1}.
\]
\end{proposition}
\begin{proof}
The entries of $A^L$ count admissible length-$L$ continuations between end states.
Summing over the terminal state corresponds to multiplying by $\mathbf{1}$, and selecting the starting state $u_6$ corresponds to left-multiplying by $e_{u_6}^{\mathsf{T}}$.
This is the matrix viewpoint recorded in Remark~\ref{rem:extension_counts_matrix_view}.
\end{proof}

\subsection{Structural comparison to the Hecke prime skeleton}
\label{subsec:protocol_hecke_comparison}

\AuditTag The Hecke prime-power recurrence (Equation~\eqref{eq:hecke_prime_power_z128}) has the same second-order skeleton as Proposition~\ref{prop:TL_prime_power_recurrence}.
In this paper this resemblance serves as a clean audit template for ``cross-scale constraint propagation'':
the refinement operators $T_L$ are explicit protocol-side objects (powers of a fixed $2\times2$ matrix),
while the Hecke operators act on modular forms and satisfy additional multiplicative relations beyond a single matrix-power semigroup (Appendix~\ref{app:hecke_prime_skeleton}).

\AuditTag Beyond the recurrence-level comparison, the same refinement-operator family also admits a rigid determinant/trace packaging:
for matrix powers $T_L=A^L$, traces $\mathrm{tr}(A^k)$ are generated by $\det(I-zA)$ via the standard log-determinant expansion.
This provides a direct bridge between refinement multiplicities and zeta/determinant bookkeeping; see Lemma~\ref{lem:rigidity_logdet_trace_series} and Remark~\ref{rem:rigidity_refinement_to_det_trace} in Section~\ref{sec:rigidity_bridge_spine}.

\AuditTag A finite enumeration check of the operator-evaluation formulas for extension multiplicities and the $\pi$-boundary subset is recorded in Table~\ref{tab:ext_boundary_operator_check} (Appendix~\ref{app:generated_tables}), generated by \texttt{scripts/exp\_ext\_boundary\_operator\_check.py}.


